\documentclass[11pt]{article}
\usepackage[margin=1in]{geometry}
\usepackage{booktabs}
\usepackage{graphicx}
\usepackage{amsmath}
\usepackage{hyperref}
\title{Scale Effects in GHAR and GNNHAR-IV Volatility Forecasting}
\author{Volatility Research Pipeline}
\date{\today}
\begin{document}
\maketitle
\section{Research Question}
We test whether graph structure becomes more valuable as the cross-sectional universe expands. For each universe, the target is 30-day realized volatility and the exogenous option-market channel is 30-day mean implied volatility. The key estimand is the out-of-sample loss reduction \(1 - L_{\mathrm{enhanced}} / L_{\mathrm{baseline}}\), reported for MSE and QLIKE.
\section{Universe-Level Results}
\begin{table}[h]
\centering
\caption{Best models and baseline losses by universe}
\begin{tabular}{lrrrrl}
\toprule
Universe & Assets & Dates & HAR QLIKE & GHAR+IV QLIKE & Best QLIKE model \\
\midrule
Dow30 & 30 & 1234 & 0.000821965 & 0.000784934 & GHAR+IV \\
S\&P100 & 99 & 1073 & 0.0033707 & 0.00323373 & GHAR+IV \\
S\&P500 & 449 & 1223 & 0.00413344 & 0.00402077 & GHAR+IV \\
\bottomrule
\end{tabular}
\end{table}
\section{Scale Gains}
\begin{table}[h]
\centering
\caption{Selected QLIKE gains by universe}
\begin{tabular}{lrrr}
\toprule
Universe & Assets & GHAR over HAR & GHAR+IV over GHAR \\
\midrule
Dow30 & 30 & 1.32\% & 3.22\% \\
S\&P100 & 99 & 0.23\% & 3.84\% \\
S\&P500 & 449 & 0.07\% & 2.66\% \\
\bottomrule
\end{tabular}
\end{table}
\section{Interpretation}
A positive GHAR-over-HAR gain means the estimated asset graph adds forecasting information beyond own-asset HAR lags. A positive GHAR+IV-over-GHAR gain means the implied-volatility channel contributes information beyond the graph. The GNN rows in the CSV output evaluate whether nonlinear message passing improves on the linear graph benchmark.
\section{Artifacts}
The machine-readable tables are \texttt{scale\_summary.csv} and \texttt{scale\_gains.csv}. The main figures are saved under \texttt{figures/}.
\end{document}
